\section{Trabalhos Relacionados}
\label{sec:relacionados}

\subsection{Propagação de atrasos e confiabilidade do tempo de viagem}

Atrasos em linhas de ônibus raramente são eventos isolados. \cite{Zhang2024} analisam atrasos de chegada em tempo real com um modelo de regressões
aparentemente não relacionadas e mostram que o atraso acumulado nas paradas anteriores é o
fator mais importante para explicar o atraso nas paradas seguintes, enquanto variáveis
meteorológicas têm papel secundário no contexto estudado. Essa ideia de acúmulo ao longo do
percurso é central para a extensão proposta neste trabalho.

Fatores da própria via também afetam a confiabilidade. \cite{KAEWUNRUEN2021100280} investigam as causas da variabilidade do tempo de viagem de
ônibus em Birmingham e apontam características do percurso, como interseções semaforizadas e
o espaçamento entre paradas, entre os elementos que mais pesam. No transporte sobre trilhos,
\cite{doi:10.1177/0361198106195500103} usam
análise computacional e simulação para aumentar a confiabilidade dos horários do S-Bahn de
Zurique, mostrando que o próprio quadro de horários pode ser uma fonte de atraso quando não
reflete os tempos reais de operação.

\subsection{Previsão com dados abertos}

A popularização do padrão GTFS e dos dados de posição em tempo real abriu espaço para modelos
de previsão de atraso. \cite{10857111} comparam modelos de
aprendizado de máquina para prever atrasos a partir de dados GTFS, e
\cite{SARHANI2024100120} discutem como combinar diferentes
fontes de dados abertos para prever atrasos no transporte ferroviário. Sistemas de
rastreamento em tempo real também são apontados como base para melhorar a eficiência da
operação e a informação ao passageiro \cite{10544098}. Esses trabalhos priorizam a acurácia da
previsão; a abordagem adotada aqui prioriza a interpretação, separando a variabilidade que se
deve ao trecho da que se deve à viagem.

\subsection{Modelos Mistos aplicados ao transporte público}

\cite{courb}, trabalho replicado aqui, analisam a linha 863 (Água
Verde) de Curitiba a partir dos dados de GPS publicados pela prefeitura, considerando um
veículo em um itinerário de 18 paradas (17 trechos). A variável resposta é o tempo de
deslocamento entre paradas consecutivas, e o agrupamento é feito pelo trecho. Os autores
comparam um modelo só com efeito fixo, um modelo com intercepto aleatório por trecho e
modelos com erro autorregressivo, incluindo um componente sazonal cujo período é o número de
trechos de uma viagem. O modelo com erro autorregressivo sazonal teve o melhor ajuste, e a
validação foi feita aplicando o modelo ajustado em uma quarta-feira aos dados da
quarta-feira seguinte.

Modelos Mistos são adequados a esse problema porque as observações têm estrutura hierárquica:
várias viagens passam pelos mesmos trechos, e cada viagem gera várias observações
\cite{pinheiro2000}. Isso permite estimar quanto da variabilidade é explicada pelo trecho e
quanto é residual, além de incorporar correlação temporal entre observações.

Este trabalho se diferencia do original em três pontos: usa outra cidade e outra fonte de
dados (a API de tempo real e o GTFS da BHTRANS); modela todas as viagens de uma linha ao longo
de 16 dias, e não um único veículo em um dia; e avalia os modelos fora da amostra com métricas
de erro, além de propor a extensão com o atraso em relação ao horário programado.
